\section{OpenAIモデルとCodex、Google Cloudへ展開}
OpenAIは、GPT-6 Astra、GPT-6.1 Sol、GPT-6 LunaとCodexのエージェント型ワークフローを、Google CloudのModel Gardenで提供すると発表した。企業はMarketplace経由で購入し、既存のクラウド利用契約に組み込めるという。まず一部企業向けの限定プレビューを開始し、一般提供は数週間以内を予定する。Google側のデータ管理の枠内で利用できる点も強調された。Microsoftが引き続き主要クラウドパートナーであることも明記されており、提供経路の拡張が企業のモデル選択肢を広げる。
\dailyxsource{OpenAI (@OpenAI)}{https://x.com/OpenAI/status/2112988118471295323}

\section{Anthropic、能動型エージェントClaude Duetを予告}
Anthropicは、Claude Opus 5.5を基盤とするアプリ「Claude Duet」の研究プレビューを発表した。利用者が席を離れている間も作業を続け、注意を要する事項を先回りして提示する一方、代理で行動する前には確認を求める設計だという。米国のPro/Max利用者を対象にウェイトリスト方式で展開し、iOS、Android、Webから利用できる。バックグラウンドでの確認作業も既存の利用枠を消費するため、利便性とコスト、判断の正確さが実運用の焦点となる。
\dailyxsource{Claude (@claudeai)}{https://x.com/claudeai/status/2112984664256241837}

\section{AIがPHPの重大な脆弱性発見を支援}
OpenAIは「Patch the Planet」の取り組みで、研究者がPHPの重大なリモートコード実行脆弱性2件を発見する際に、同社のモデルが貢献したと報告した。PHP 5.4以降に関係し、広範なWebサイトに影響し得るという。約2か月の協調的な開示を経て修正版が公開されたとされる。コミュニティノートでは独立した発見者の関与も指摘されており、発見の帰属や影響規模は単純化できない。具体的なCVE番号や悪用の有無はGrokの収集時点で未確定として扱う。
\dailyxsource{OpenAI (@OpenAI)}{https://x.com/OpenAI/status/2112907856018309586}

\section{LiteLLMの配布版に認証情報窃取マルウェア}
AIモデルへの接続に使われるPythonパッケージLiteLLMのv1.107.0について、PyPIとDockerの配布物に悪意あるコードが混入したと報告された。Socketの早期投稿とLiteLLMの公式説明によれば、クラウドやSSH、AIサービスの認証情報などを収集する機能が含まれていた。侵害された配布物は引き下げられ、公式案内では1.106.2が安全な版として示されている。企業向けプライベートリリースは影響を受けないという。AI開発の共通部品を経由した供給網攻撃として、利用版の確認と認証情報の保護が重要となる。
\dailyxsource{Socket (@SocketSecurity)、窓内の早期報告}{https://x.com/SocketSecurity/status/2112907732877775016}

\section{Gemini 4 Argon、人体解剖の3Dデモを披露}
Google DeepMindは、Gemini 4 Argonが3つのオープンソースデータセットを組み合わせて作成したとする、人体解剖のインタラクティブな3D可視化を紹介した。筋肉などのレイヤー切り替え、個別部位の分離、心臓内部へのズームが可能なデモである。異なるデータを統合し、操作可能な可視化に変えるモデルの能力を示す例として注目された。ただし、使用データセットの詳細や再現方法を説明する技術資料は、今回のGrok収集では確認されていない。
\dailyxsource{Google DeepMind (@GoogleDeepMind)}{https://x.com/GoogleDeepMind/status/2112974723659247896}
